% !TeX root = ../../book.tex
\section{弱归纳法}
\secbegin{secWeakInduction}
\index{induction!weak|(}

正如递归利用自然数的结构来\textit{定义}涉及自然数的\textit{表达式}一样，归纳法也利用相同的结构来\textit{证明}有关自然数的\textit{结果}。

\begin{restatable}[弱归纳原理]{theorem}{rthmWeakInduction}
\label{thmWeakInduction}
\index{induction!on $\mathbb{N}$ (weak)}
\index{weak induction principle}
设 $p(n)$ 为带有自由变量 $n \in \mathbb{N}$ 的逻辑公式，并设 $n_0 \in \mathbb{N}$。如果
\begin{enumerate}[(i)] 
\item $p(n_0)$ 为真；
\item 对于所有 $n \ge n_0$，如果 $p(n)$ 为真，则 $p(n+1)$ 为真；
\end{enumerate}
那么 $p(n)$ 对于所有 $n \ge n_0$ 为真。
\end{restatable}

\begin{cproof}
定义 $X = \{ n \in \mathbb{N} \mid p(n_0 + n) \text{ 为真} \}$；也就是说，给定自然数 $n$，我们有 $n \in X$ 当且仅当 $p(n_0 + n)$ 为真。则
\begin{itemize}
\item $0 \in X$，因为 $n_0 + 0 = n_0$ 且根据 (i) $p(n_0)$ 为真。
\item 设 $n \in \mathbb{N}$ 并假设 $n \in X$。则 $p(n_0+n)$ 为真，因为 $n_0 + n \ge n_0$ 且 $p(n_0+n)$ 为真，根据(ii) 我们有 $p(n_0+n+1)$ 为真。所以 $n+1 \in X$。
\end{itemize}

因此，根据 \Cref{defNotionOfNaturalNumbers}(iii) 我们有 $\mathbb{N} \subseteq X$。所以，$p(n_0 + n)$ 对于所有 $n \in \mathbb{N}$ 都为真。而这相当于说 $p(n)$ 对于所有 $n \ge n_0$ 都为真。
\end{cproof}

\begin{strategy}[{(弱)} 归纳法证明]
\index{proof!by weak induction}
要想证明 $\forall n \ge n_0,\, p(n)$ 形式的命题，只需证明：
\begin{itemize}
\item $p(n_0)$ 为真；
\item 对于所有 $n \ge n_0$，如果 $p(n)$ 为真，则 $p(n+1)$ 为真。
\end{itemize}
归纳法证明引入了一些术语：
\begin{itemize}
\item $p(n_0)$ 的证明称为\textbf{基本情况}；
\item $\forall n \ge n_0,\, (p(n) \Rightarrow p(n+1))$ 的证明称为\textbf{归纳步骤}；
\item 在归纳步骤中，假设 $p(n)$ 为真称为\textbf{归纳假设}；
\item 在归纳步骤中，命题 $p(n+1)$ 称为\textbf{归纳结果}。
\end{itemize}
\end{strategy}

下图展示了弱归纳原理。

\begin{center}
\begin{tikzpicture}
\draw[fill=indstepcol] (6,-1) rectangle (8,1);
\draw (0,0) node[diamond, draw, text width=27, text centered, fill=indbasecol] {$n_0$};
\draw (2,0) node[circle, draw, text width=27, text centered, fill=white] {$n_0+1$};
\draw (3.5,0) node {$\cdots$};
\draw (5,0) node[circle, draw, text width=27, text centered, fill=white] {$n-1$};
\draw (7,0) node[circle, draw, text width=27, text centered, fill=white] {$n$};
\draw (9.5,0) node[circle, draw, dashed, text width=27, text centered, fill=white](nplusone) {$n+1$};
\draw[->] (8,0) -- (nplusone);
\end{tikzpicture}
\end{center}

对上图做一些解释：
\begin{itemize}
\item 阴影菱形代表基本情况 $p(n_0)$；
\item 正方形代表归纳假设 $p(n)$；
\item 虚线圆形代表归纳目标 $p(n+1)$；
\item 箭头代表我们在归纳步骤中要证明的蕴涵关系。
\end{itemize}

我们将使用类似的图来描述本节中的其他归纳原理。

\begin{proposition}
\label{propSumConsecInduction}
设 $n \in \mathbb{N}$。则 $\displaystyle\sum_{k=1}^n k = \frac{n(n+1)}{2}$
\end{proposition}

\begin{cproof}
我们在 $n \ge 0$ 上用归纳法进行证明。

\begin{itemize}
\item (\textbf{基本情况})我们需要证明 $\displaystyle\sum_{k=1}^0 k = \frac{0(0+1)}{2}$。

这明显成立。因为 $\dfrac{0(0+1)}{2} = 0$，而根据\Cref{defSumOfRealNumbers} $\displaystyle\sum_{k=1}^0 k = 0$。

\item (\textbf{归纳步骤}) 设 $n \ge 0$ 并假设 $\displaystyle\sum_{k=1}^n k = \dfrac{n(n+1)}{2}$；此为归纳假设。

我们需要证明 $\displaystyle\sum_{k=1}^{n+1} k = \frac{(n+1)(n+2)}{2}$；此为归纳结果。

通过计算得：
\begin{align*}
\sum_{k=1}^{n+1} k &= \left( \sum_{k=1}^n k \right) + (n+1) && \text{根据 \Cref{defSumOfRealNumbers}} \\
&= \frac{n(n+1)}{2} + (n+1) && \text{根据归纳假设}\\
&= (n+1)\left( \frac{n}{2} + 1 \right) && \text{提取公因式} \\
&= \frac{(n+1)(n+2)}{2} && \text{整理}
\end{align*}
\end{itemize}
由归纳得出结果。
\end{cproof}

在继续深入之前，让我们回顾一下 \Cref{propSumConsecInduction} 的证明，强调一些通过归纳法编写证明的有效方法。
\begin{itemize}
\item 我们在证明开始时指出这是一个归纳证明。虽然本节清楚表明大多数证明将通过归纳法进行证明，但情况并非总是如此，因此最好指出当下的证明策略。
\item 证明中清楚地标记出基本情况和归纳步骤。从数学角度来看，这并不是严格\textit{必须的}，但它可以帮助读者了解证明结构并确定每一步的目标是什么。
\item 我们通过写下`设 $n \ge n_0$ 并假设 [\dots{}\textit{此为归纳假设}\dots{}]'开始归纳步骤。归纳步骤通常都是这样开始的，因为在归纳步骤中证明的命题的形式为 $\forall n \ge n_0,\, (p(n) \Rightarrow \cdots)$。
\item 在证明基本情况或归纳步骤之前，我们写下该部分试图证明的内容。这很有用，因为它有助于提醒我们（以及阅读证明的人）我们的目标是什么。
\end{itemize}

下一个命题也是通过 $n \ge 0$ 归纳得出的，在其证明过程中注意这些特征。

\begin{proposition}
对于所有 $n \in \mathbb{N}$，自然数 $n^3-n$ 可被 $3$ 整除。
\end{proposition}

\begin{cproof}
我们在 $n \ge 0$ 上用归纳法进行证明。
\begin{itemize}
\item (\textbf{基本情况}) 我们需要证明 $0^3-0$ 能被 $3$ 整除。
\[ 0^3-0 = 0 = 3 \times 0 \]
所以 $0^3-0$ 能被 $3$ 整除。

\item (\textbf{归纳步骤}) 设 $n \in \mathbb{N}$ 并假设 $n^3-n$ 能被 $3$ 整除。则可以写成 $n^3-n = 3k$, $k \in \mathbb{Z}$ 的形式。

我们需要证明 $(n+1)^3 - (n+1)$ 可以被 $3$ 整除；换句话说，我们需要找到某自然数 $\ell$ 使得
\[ (n+1)^3 - (n+1) = 3\ell \]
通过计算得
\begin{align*}
& (n+1)^3 - (n+1) && \\
&= (n^3 + 3n^2 + 3n + 1) - n - 1 && \text{展开}\\
&= n^3 - n + 3n^2 + 3n + 1 - 1 && \text{整理} \\
&= n^3 - n + 3n^2 + 3n && \text{因为 $1 - 1 = 0$} \\
&= 3k + 3n^2 + 3n && \text{根据递归假设} \\
&= 3(k+n^2+n) && \text{提取公因数}
\end{align*}
因此，对于某些 $\ell \in \mathbb{Z}$，我们将 $(n+1)^3-(n+1)$ 表示为 $3\ell$ 形式；具体来说，$\ell=k+n^2+n$。
\end{itemize}

由归纳得出结果。
\end{cproof}

\begin{exercise}
\label{exSumOfPowersOf2}
用归纳法证明对于所有 $n \in \mathbb{N}$, $\displaystyle \sum_{k=0}^n 2^k = 2^{n+1} - 1$。
\end{exercise}

下面的命题可以用归纳法证明，但基本情况不为零。

\begin{proposition}
\label{propWeakInductionNonzeroBaseCaseExample}
对于所有 $n \ge 4$，我们有 $3n < 2^n$。
\end{proposition}

\begin{cproof}
我们在 $n \ge 4$ 上用归纳法进行证明。
\begin{itemize}
\item (\textbf{基本情况}) $p(4)$ 代表陈述 $3 \cdot 4 < 2^4$。这显然成立，因为 $12 < 16$。

\item (\textbf{归纳步骤}) 假设 $n \ge 4$ 且 $3n < 2^n$。我们要证明 $3(n+1) < 2^{n+1}$。
\begin{align*}
3(n+1) &= 3n+3 && \text{展开} \\
&< 2^n + 3   && \text{根据归纳假设} \\
&< 2^n + 2^4 && \text{因为 $3 < 16 = 2^4$} \\
&\le 2^n + 2^n && \text{因为 $n \ge 4$} \\
&= 2 \cdot 2^n && \text{化简} \\
&= 2^{n+1} && \text{化简}
\end{align*}
因此我们证明出 $3(n+1) < 2^{n+1}$。
\end{itemize}
由归纳得出结果。
\end{cproof}

请注意，\Cref{propWeakInductionNonzeroBaseCaseExample} 中的证明没有说明 $p(n)$ 对于 $n=0,1,2,3$ 的真假。为了断言这些案例是错误的，您需要单独呈现它们：
\begin{itemize}
\item $3 \times 0 = 0$ 且 $2^0 = 1$，因此 $p(0)$ 为真；
\item $3 \times 1 = 3$ 且 $2^1 = 2$，因此 $p(1)$ 为假；
\item $3 \times 2 = 6$ 且 $2^2 = 4$，因此 $p(2)$ 为假；
\item $3 \times 3 = 9$ 且 $2^3 = 8$，因此 $p(3)$ 为假。
\end{itemize}
因此我们推断，当 $n=0$ 或 $n \ge 4$ 时，$p(n)$ 为真，当 $n \in \{ 1, 2, 3 \}$ 时，$p(n)$ 为假。

有时，用归纳法`证明'可能看起来完全是无稽之谈。以下是`归纳失败'的经典示例：

\begin{example}
\label{exHorseInductionFail}
下面的论证可以证明每匹马都是相同颜色。
\begin{itemize}
\item (\textbf{基本情况}) 假设只有一匹马。这匹马跟它自己具有相同的颜色，因此基本情况显然成立。
\item (\textbf{归纳步骤}) 假设 $n$ 匹马的每个集合都是相同颜色。设 $X$ 为 $n+1$ 匹马的集合。去掉 $X$ 中的第一匹马，根据归纳假设，我们得到后 $n$ 匹马的颜色相同。从 $X$ 中删除最后一匹马，我们得到前 $n$ 匹马的颜色相同。因此 $X$ 中所有马都是相同颜色。
\end{itemize}
根据归纳法，证明完毕。
\end{example}

\begin{exercise}
写出 \Cref{exHorseInductionFail} 尝试对所有 $n \ge 1$ 证明的语句 $p(n)$。证明基本情况的证明是正确的，然后用量词准确写出归纳步骤要证明的命题。解释为什么归纳步骤中的论证未能证明这个命题。
\end{exercise}

有几种方法可以避免类似 \Cref{exHorseInductionFail} 的情况，只需在编写证明时多加考虑即可。下面是一些提示：

\begin{itemize}
\item 明确陈述 $p(n)$。在`所有马都是相同颜色'的陈述中，并不清楚归纳变量到底是什么。然而，我们可以说：
\begin{center}
设 $p(n)$ 为陈述`$n$ 匹马的每个集合都具有相同颜色'。
\end{center}

\item 在归纳步骤中明确参考基本情况 $n_0$。在 \Cref{exHorseInductionFail} 中，我们的归纳假设只是简单地表述为`假设 $n$ 匹马的每个集合具有相同的颜色'。我们可以换种说法：
\begin{center}
设 $n \ge 1$ 并假设 $n$ 匹马的每个集合具有相同颜色。
\end{center}
我们可能已经发现了接下来发生的错误。
\end{itemize}

下面是一些弱归纳法证明的例子。

\begin{proposition}
对于所有 $n \in \mathbb{N}$，我们有 $\displaystyle \sum_{k=0}^n k^3 = \left( \sum_{k=0}^n k \right)^2$。
\end{proposition}

\begin{cproof}
在 \Cref{propWeakInductionNonzeroBaseCaseExample} 中我们证明了对于所有 $n \in \mathbb{N}$, $\displaystyle\sum_{k=0}^n k = \frac{n(n+1)}{2}$，因此只需证明对于所有 $n \in \mathbb{N}$, 
\[ \sum_{k=0}^n k^3 = \frac{n^2(n+1)^2}{4} \]

我们在 $n \ge 0$ 上用归纳法进行证明。
\begin{itemize}
\item (\textbf{基本情况}) 我们需要证明 $\displaystyle 0^3 = \frac{0^2(0+1)^2}{4}$。这明显成立，因为公式两边都为 $0$。
\item (\textbf{归纳步骤}) 给定 $n \in \mathbb{N}$ 并假设 $\displaystyle \sum_{k=0}^n k^3 = \frac{n^2(n+1)^2}{4}$。我们需要证明 $\displaystyle \sum_{k=0}^{n+1} k^3 = \frac{(n+1)^2(n+2)^2}{4}$。这也成立，因为:
\begin{align*}
\sum_{i=0}^{n+1} k^3 & = \sum_{i=0}^n k^3 + (n+1)^3 && \text{根据加法定义} \\
&= \frac{n^2(n+1)^2}{4} + (n+1)^3 && \text{根据归纳假设} \\
&= \frac{n^2(n+1)^2 + 4(n+1)^3}{4} && \text{(\textit{通分})} \\
&= \frac{(n+1)^2(n^2 + 4(n+1))}{4} && \text{(\textit{提取公因式})} \\
&= \frac{(n+1)^2(n+2)^2}{4} && \text{(\textit{完全平方公式})}
\end{align*}
\end{itemize}
通过归纳法得出结果。
\end{cproof}

下一命题中，我们证明了众所周知的实数\textit{算术级数}求和公式的正确性。

\begin{proposition}
设 $a,d \in \mathbb{R}$。则对于所有 $n \in \mathbb{N}$, 
\[ \sum_{k=0}^n (a+kd) = \frac{(n+1)(2a+nd)}{2} \]
\end{proposition}

\begin{cproof}
我们在 $n \ge 0$ 上用归纳法进行证明。
\begin{itemize}
\item (\textbf{基本情况}) 我们需要证明 $\displaystyle\sum_{k=0}^0 (a+kd) = \frac{(0+1)(2a+0d)}{2}$。这明显成立，因为
\[ \sum_{k=0}^0 (a+kd) = a + 0d = a = \frac{2a}{2} = \frac{1 \cdot (2a)}{2} = \frac{(0+1)(2a+0d)}{2} \]

\item (\textbf{归纳步骤}) 给定 $n \in \mathbb{N}$ 并假设 $\displaystyle\sum_{k=0}^n (a+kd) = \frac{(n+1)(2a+nd)}{2}$。我们需要证明：
\[ \sum_{k=0}^{n+1} (a+kd) = \frac{(n+2)(2a+(n+1)d)}{2} \]
这也成立，因为：
\begin{align*}
&\sum_{k=0}^{n+1} (a+kd) && \\
&= \sum_{k=0}^n (a+kd) + (a+(n+1)d) && \text{根据加法定义} \\
&= \frac{(n+1)(2a+nd)}{2} + (a + (n+1)d) && \text{根据归纳假设} \\
&= \frac{(n+1)(2a+nd) + 2a + 2(n+1)d}{2} && \text{(\textit{通分})} \\
&= \frac{(n+1) \cdot 2a + (n+1) \cdot nd + 2a + 2(n+1)d}{2} && \text{(\textit{乘法展开})} \\
&= \frac{2a(n+1+1) + (n+1)(nd + 2d)}{2} && \text{(\textit{提取公因式})} \\
&= \frac{2a(n+2) + (n+1)(n+2)d}{2} && \text{(\textit{整理})} \\
&= \frac{(n+2)(2a + (n+1)d)}{2} && \text{(\textit{提取公因式})}
\end{align*}
\end{itemize}
通过归纳法得出结果。
\end{cproof}

以下练习泛化了 \Cref{exSumOfPowersOf2} 以证明实数\textit{几何级数}求和公式的正确性。

\begin{exercise}
\label{exFormulaForGeometricProgression}
设 $a,r \in \mathbb{R}$ 其中 $r \ne 1$。则对于所有 $n \in \mathbb{N}$, 
\[ \sum_{k=0}^n ar^k = \frac{a(1-r^{n+1})}{1-r} \]
\end{exercise}

到目前为止，在本节中，我们通过归纳法证明的大部分结论都只涉及数字，特别是涉及自然数的恒等式。但归纳法证明的范围并不限于这些结论 --- 例如，下一个练习涉及集合和双射。

\begin{exercise}
设 $n \in \mathbb{N}$ 并设 $\{ X_k \mid 1 \le k \le n \}$ 为集合族。在 $n$ 上用归纳法证明 $\displaystyle \prod_{k=1}^{n+1} X_k \to \left( \prod_{k=1}^n X_k \right) \times X_n$ 为双射。
\hintlabel{exProductOfSuccNSets}{%
要定义双射，请考虑两个集合的元素是什么样的：$\prod_{k=1}^{n+1} X_k$ 中的元素类似 $(a_1, a_2, \dots, a_n, a_{n+1})$，其中对于每个 $1 \le k \le n+1$, $a_k \in X_k$。另一方面，$\left( \prod_{k=1}^n X_k \right) \times X_{n+1}$ 中的元素类似 $((a_1, a_2, \dots, a_n), a_{n+1})$。
}
\end{exercise}

\subsection*{二项式和阶乘}

归纳法对证明二项式系数 $\binom{n}{k}$ 和阶乘 $n!$ 非常有用。

\begin{example}
\label{exSumOfBinomialCoefficients}
在 $n$ 上，用归纳法证明 $\sum_{i=0}^n \binom{n}{i} = 2^n$。
\begin{itemize}
\item (\textbf{基本情况}) 我们要证明 $\binom{0}{0} = 1$ 且 $2^0=1$。根据二项式系数和指数的定义，这两部分都成立。
\item (\textbf{归纳步骤}) 给定 $n \ge 0$ 并假设
\[ \sum_{i=0}^n \binom{n}{i} = 2^n \]
我们需要证明
\[ \sum_{i=0}^{n+1} \binom{n+1}{i} = 2^{n+1} \]
这是成立的，因为

\begin{align*}
&\sum_{i=0}^{n+1} \binom{n+1}{i} && \\
&= \binom{n+1}{0} + \sum_{i=1}^{n+1} \binom{n+1}{i} && \text{拆分求和公式} \\
&= 1 + \sum_{j=0}^n \binom{n+1}{j+1} && \text{设 $j=i-1$} \\
&= 1 + \sum_{j=0}^n \left(\binom{n}{j} + \binom{n}{j+1}\right) && \text{根据 \Cref{defBinomialCoefficientRecursive}} \\
&= 1 + \sum_{j=0}^n \binom{n}{j} + \sum_{j=0}^n \binom{n}{j+1} && \text{分别求和}
\end{align*}

根据递归假设 $\sum_{j=0}^n \binom{n}{j} = 2^n$。此外，使用 $k=j+1$ 重新索引求和会产生
\[ \sum_{j=0}^n \binom{n}{j+1} = \sum_{k=1}^{n+1} \binom{n}{k} = \sum_{k=1}^n \binom{n}{k} + \binom{n}{n+1} \]
根据递归假设，我们有
\[ \sum_{k=1}^n \binom{n}{k} = \sum_{k=0}^n \binom{n}{k} - \binom{n}{0} = 2^n-1 \]
且 $\binom{n}{n+1} = 0$，所以 $\sum_{j=0}^n \binom{n}{j+1} = 2^n-1$.

综上可得
\begin{align*}
1 + \sum_{j=0}^n \binom{n}{j} + \sum_{j=0}^n \binom{n}{j+1}
&= 1 + 2^n + (2^n-1) \\
&= 2 \cdot 2^n \\
&= 2^{n+1}
\end{align*}
至此归纳步骤结束。
\end{itemize}
通过归纳法得出结果。
\end{example}

\begin{theorem}
\label{thmBinomAsFactorialByInduction}
设 $n,k \in \mathbb{N}$。则
\[ \binom{n}{k} =
\begin{cases}
\dfrac{n!}{k!(n-k)!} & \text{如果 } k \le n \\
0 & \text{如果 } k > n
\end{cases} \]
\end{theorem}
\begin{cproof}
我们对 $n$ 进行归纳。
\begin{itemize}
\item (\textbf{基本情况}) 当 $n=0$ 时，我们需要证明对于所有 $k \le 0$, $\binom{0}{k} = \frac{0!}{k!(-k)!}$，且对于所有 $k>0$, $\binom{0}{k} = 0$。

如果 $k \le 0$ 则 $k=0$，因为 $k \in \mathbb{N}$。因此我们需要证明
\[ \binom{0}{0} = \frac{0!}{0!0!} \]
这显然成立，因为 $\binom{0}{0}=1$ 且 $\frac{0!}{0!0!} = \frac{1}{1 \times 1} = 1$。

如果 $k>0$ 则根据\Cref{defBinomialCoefficientRecursive} $\binom{0}{k} = 0$。
\item (\textbf{归纳步骤}) 给定 $n \in \mathbb{N}$ 并假设对于所有 $k \le n$, $\binom{n}{k} = \frac{n!}{k!(n-k)!}$ 且对于所有 $k > n$, $\binom{n}{k} = 0$。

我们需要证明对于所有 $k \le n+1$，我们有
\[ \binom{n+1}{k} = \frac{(n+1)!}{k!(n+1-k)!} \]
且对于所有 $k > n+1$, $\binom{n+1}{k} = 0$。

所以给定 $k \in \mathbb{N}$。存在 4 种可能的情况：要么 (i) $k=0$，要么 (ii) $0 < k \le n$，要么 (iii) $k=n+1$，要么 (iv) $k>n+1$。情况 (i), (ii), (iii) 中，我们需要证明 $\binom{n+1}{k}$ 的阶乘公式；情况 (iv) 中，我们需要证明 $\binom{n+1}{k} = 0$。

\begin{enumerate}[(i)]
\item 假设 $k=0$。则根据\Cref{defBinomialCoefficientRecursive}可得 $\binom{n+1}{0} = 1$，且
\[ \frac{(n+1)!}{k!(n+1-k)!} = \frac{(n+1)!}{0!(n+1)!} = 1 \]
应为 $0!=1$。所以 $\binom{n+1}{0} = \frac{(n+1)!}{0!(n+1)!}$。
\item 如果 $0 < k \le n$ 则对于某自然数 $\ell < n$, $k=\ell+1$。那么 $\ell+1 \le n$，因此我们可以使用归纳假设将阶乘公式应用于 $\binom{n}{\ell}$ 和 $\binom{n}{\ell+1}$。所以
\begin{align*}
&\binom{n+1}{k} && \\
&= \binom{n+1}{\ell+1} && \text{因为 $k=\ell+1$} \\
&= \binom{n}{\ell} + \binom{n}{\ell+1} && \text{根据 \Cref{defBinomialCoefficientRecursive}} \\
&= \frac{n!}{\ell!(n-\ell)!} + \frac{n!}{(\ell+1)!(n-\ell-1)!} && \text{根据递归假设}
\end{align*}
请注意
\[ \frac{n!}{\ell!(n-\ell)!} = \frac{n!}{\ell!(n-\ell)!} \cdot \frac{\ell+1}{\ell+1} = \frac{n!}{(\ell+1)!(n-\ell)!} \cdot (\ell+1) \]
且
\begin{center}
$\dfrac{n!}{(\ell+1)!(n-\ell-1)!} = \dfrac{n!}{(\ell+1)!(n-\ell-1)!} \cdot \dfrac{n-\ell}{n-\ell} = \dfrac{n!}{(\ell+1)!(n-\ell)!} \cdot (n-\ell)$
\end{center}

综上可得
\begin{align*}
&\frac{n!}{\ell!(n-\ell)!} + \frac{n!}{(\ell+1)!(n-\ell-1)!} && \\
&= \frac{n!}{(\ell+1)!(n-\ell)!} \cdot [(\ell+1) + (n-\ell)] \\
&= \frac{n!(n+1)}{(\ell+1)!(n-\ell)!} && \\
&= \frac{(n+1)!}{(\ell+1)!(n-\ell)!}
\end{align*}
因此 $\binom{n+1}{\ell+1} = \frac{(n+1)!}{(\ell+1)!(n-\ell)!}$。搞定；的确
\[ \frac{(n+1)!}{(\ell+1)!(n-\ell)!} = \frac{(n+1)!}{k!(n+1-k)!} \]
因为 $k=\ell+1$。
\item 如果 $k=n+1$，则
\begin{align*}
\binom{n+1}{k} &= \binom{n+1}{n+1} && \text{因为 $k=n+1$} \\
&= \binom{n}{n} + \binom{n}{n+1} && \text{根据 \Cref{defBinomialCoefficientRecursive}} \\
&= \frac{n!}{n!0!} + 0 && \text{根据递归假设} \\
&= 1
\end{align*}
且 $\frac{(n+1)!}{(n+1)!0!}=1$，这两个量再次相等。
\item 如果 $k>n+1$，则对于某个 $\ell > n$, $k=\ell+1$，因此根据 \Cref{defBinomialCoefficientRecursive} 和递归假设可得
\[ \binom{n+1}{k} = \binom{n+1}{\ell+1} \overset{\text{IH}}{=} \binom{n}{\ell}+\binom{n}{\ell+1} = 0+0 = 0 \]

\end{enumerate}
\end{itemize}
\end{cproof}

初读时，这个证明又臭又长，特别是在归纳步骤中，我们需要分为四种情况。我们将在 \Cref{secCountingPrinciples} 中给出一个更简单的证明（参见 \Cref{thmBinomAsFactorial}），其中我们通过将两个集合的元素一一对应来\textit{组合地}证明该陈述。

我们可以使用 \Cref{thmBinomAsFactorialByInduction} 来证明很多涉及二项式系数的非常有用的恒等式。

\begin{example}
设 $n,k,\ell \in \mathbb{N}$ 其中 $\ell \le k \le n$ 则 \[ \binom{n}{k}\binom{k}{\ell} = \binom{n}{\ell}\binom{n-\ell}{k-\ell} \]
证明：
\begin{align*}
& \binom{n}{k} \binom{k}{\ell} && \\
&= \frac{n!}{k!(n-k)!} \cdot \frac{k!}{\ell!(k-\ell!)} && \text{根据 \Cref{thmBinomAsFactorialByInduction}} \\
&= \frac{n!k!}{k!\ell!(n-k)!(k-\ell)!} && \text{合并} \\
&= \frac{n!}{\ell!(n-k)!(k-\ell)!} && \text{消掉 } k!\\
&= \frac{n!(n-\ell)!}{\ell!(n-k)!(k-\ell)!(n-\ell)!} && \text{乘以 } \frac{(n-\ell)!}{(n-\ell)!} \\
&= \frac{n!}{\ell!(n-\ell!)} \cdot \frac{(n-\ell)!}{(k-\ell)!(n-k)!} && \text{分离} \\
&= \frac{n!}{\ell!(n-\ell!)} \cdot \frac{(n-\ell)!}{(k-\ell)!((n-\ell)-(k-\ell))!} && \text{整理} \\
&= \binom{n}{\ell} \binom{n-\ell}{k-\ell} && \text{根据 \Cref{thmBinomAsFactorialByInduction}}
\end{align*}
\end{example}

\begin{exercise}
证明对于所有 $n,k \in \mathbb{N}$, $\dbinom{n}{k} = \dbinom{n}{n-k}$，其中 $k \le n$。
\end{exercise}

二项式系数在初等代数中的一个非常有用的应用是二项式定理。

\begin{theorem}[二项式定理]
\label{thmBinomialTheorem}
设 $n \in \mathbb{N}$ 且 $x,y \in \mathbb{R}$。则
\[ (x+y)^n = \sum_{k=0}^n \binom{n}{k} x^k y^{n-k} \]
\end{theorem}
\begin{cproof}
当 $y=0$ 时，我们有对于所有 $k<n$, $y^{n-k}=0$，所以方程简化为
\[ x^n = x^ny^{n-n} \]
这显然成立，因为 $y^0=1$。因此，对于其余的证明，我们假设 $y \ne 0$。

现在我们先简化为 $y=1$ 时的情况；然后扩展到任意 $y \ne 0$。

对 $n$ 用归纳法证明 $(1+x)^n = \sum_{k=0}^n \binom{n}{k} x^k$。
\begin{itemize}
\item (\textbf{基本情况}) $(1+x)^0=1$ 且 $\binom{0}{0} x^0 = 1 \cdot 1 = 1$，所以当 $n=0$ 时成立。
\item (\textbf{归纳步骤}) 给定 $n \in \mathbb{N}$ 并假设
\[ (1+x)^n = \sum_{k=0}^n \binom{n}{k} x^k \]

我们需要证明 $(1+x)^{n+1} = \sum_{k=0}^{n+1} \binom{n+1}{k} x^k$。

\begin{align*} \hspace{-30pt}
& (1+x)^{n+1} && \\
&= (1+x)(1+x)^n && \text{根据指数定律} \\
&= (1+x) \cdot \sum_{k=0}^n \binom{n}{k} x^k && \text{根据递归假设} \\
&= \sum_{k=0}^n \binom{n}{k} x^k + x \cdot \sum_{k=0}^n \binom{n}{k} x^{k}&& \text{展开 $(x+1)$} \\
&= \sum_{k=0}^n \binom{n}{k} x^k + \sum_{k=0}^n \binom{n}{k} x^{k+1} && \text{分配 $x$} \\
&= \sum_{k=0}^n \binom{n}{k} x^k + \sum_{k=1}^{n+1} \binom{n}{k-1} x^k && \text{后一个求和 $k \to k-1$} \\
&= \binom{n}{0}x^0 + \sum_{k=1}^n \left( \binom{n}{k} + \binom{n}{k-1} \right) x^k + \binom{n}{n}x^{n+1} && \text{将求和拆开} \\
&= \binom{n}{0}x^0 + \sum_{k=1}^n \binom{n+1}{k} x^k + \binom{n}{n}x^{n+1} && \text{根据 \Cref{defBinomialCoefficientRecursive}} \\
&= \binom{n+1}{0}x^0 + \sum_{k=1}^n \binom{n+1}{k} x^k + \binom{n+1}{n+1}x^{n+1} && \text{参见下方 $(*) 部分$ below} \\
&= \sum_{k=0}^{n+1} \binom{n+1}{k} x^k &&
\end{align*}
$(*)$ 标步骤成立是因为
\[ \binom{n}{0} = 1 = \binom{n+1}{0} \qquad \text{and} \qquad \binom{n}{n} = 1 = \binom{n+1}{n+1} \]
\end{itemize}
通过归纳法，我们证明了对于所有 $n \in \mathbb{N}$, $(1+x)^n = \sum_{i=0}^n \binom{n}{k} x^k$。

当 $y \ne 0$ 且不一定等于 $1$ 时，我们有
\[ (x+y)^n = y^n \cdot \left( 1 + \frac{x}{y} \right)^n = y^n \cdot \sum_{k=0}^n \binom{n}{k} \left( \frac{x}{y} \right)^k = \sum_{k=0}^n \binom{n}{k} x^k y^{n-k} \]
中间的方程遵循我们刚刚证明的；最左边和最右边的方程是简单的代数整理。
\end{cproof}

\begin{example}
在 \Cref{exSumOfBinomialCoefficients} 中我们看到
\[ \sum_{k=0}^n \binom{n}{k} = 2^n \]
这可以从二项式定理中快速推出，因为
\[ 2^n = (1+1)^n = \sum_{k=0}^n \binom{n}{k} \cdot 1^k \cdot 1^{n-k} = \sum_{k=0}^n \binom{n}{k} \]
同理，在 \Cref{exAlternatingSumOfBinomialCoefficientsIsZero} 中，我们证明了二项式系数的交替和为零；也就是说，对于 $n \in \mathbb{N}$，我们有
\[ \sum_{k=0}^n (-1)^k \binom{n}{k} = 0 \]
通过使用二项式定理大大简化了证明。事实上，根据二项式定理，我们有
\[ 0 = 0^n = (-1+1)^n = \sum_{k=0}^n \binom{n}{k}(-1)^k1^{n-k} = \sum_{k=0}^n (-1)^k \binom{n}{k} \]
使用\textit{枚举组合}可以更优雅、更快速、更轻松地证明这两个恒等式。这将是 \Cref{secCountingPrinciples} 中涵盖的主题。
\end{example}

\begin{exercise}
利用二项式定理证明
\[ \sum_{i=0}^n (-1)^i \binom{n}{i} = 0 \]
\hintlabel{exAlternatingSumOfBinomialCoefficientsIsZero}{%
注意 $(-1)^i \dbinom{n}{i} = \dbinom{n}{i} \cdot (-1)^i \cdot 1^{n-i}$。
}
\end{exercise}

\index{induction!weak|)}